\section{Conclusion}
\label{sec:conclusion}

In this pilot, text instructions alone did not prevent image-editing APIs
from changing a predefined facial keep zone. A landmark-derived mask,
alignment, and client-side feathered composite reduced that measured off-target
change while retaining similar target-region pixel change, uniformly across the
six editors tested and at a median provider cost of \$\MedCost{} per image. The
tested Qwen inpainting model did not outperform that post-hoc control, but
one model cannot support a general conclusion about masked inpainting.

The study does not validate cosmetic-surgery outcome prediction, and
client-side compositing should be read as a control for off-target pixel
changes, not as evidence that a preview is clinically accurate or safe for
patient counseling. The main analysis contains \NumPrimaryFaces{} faces, one
generated output per cell, no surgeon ratings, limited profile support, and
only an identity-embedding measure of movement toward the postoperative
photographs; by that measure the edits moved slightly away from the outcome,
not closer. The next study should add repeated generations, pose-robust
whole-image localization, and anatomy-aware postoperative accuracy metrics. It will
also need multiple blinded surgeons with inter-rater agreement, a larger and
demographically characterized cohort, and immutable or self-hosted model
baselines. What the pilot does establish is
narrower: when an editor exposes no internals, region confinement can still be
enforced entirely from the client.
